\documentclass{article}

\usepackage[margin=1in]{geometry}
\usepackage{amsmath, amssymb}
\usepackage{array}
\usepackage{xcolor}

\renewcommand{\arraystretch}{1.4}
\setlength{\tabcolsep}{0pt}

\title{Working Document -- Conversational Advertisement}
\author{}
\date{}

\begin{document}

\maketitle

\section{Goal of this Document}

% PROBLEM SETTING
We aim to identify and model meaningful advertising opportunities within multi‑turn conversational AI agents, and to benchmark different advertising strategies, conversational recommendation strategies, and ad‑intensity policies. Our work leverages user psychological traits (OCEAN), multi‑turn conversation history, and user history to understand when and how advertisements should be introduced.
\\ 

% OPTIMIZATION TARGET
The main goal is to quantify the trade‑offs between user trust (which is related to but distinct from user value), advertiser trust, and sustainable platform revenue, going beyond simple metrics such as CTR. These are the exact trade‑offs that many LLM providers are urgently trying to figure out.
\\

% VALIDATION LAYER
To validate and refine these insights, we plan a user study that examines the neurophysiological effects of different ad‑intensity and persuasion regimes, using EEG and eye‑tracking data to identify the “sweet spots” where user trust, advertiser welfare, and platform sustainable revenue are jointly optimized.

% --------------------------------------------------
\section{Research Questions}

\begin{itemize}
    \item \textbf{RQ1:} How is the ad format affecting ad conversion?
    \item \textbf{RQ2:} How is the ad format affecting the subjective user experience, including perceived trust in the AI agent?
    \item \textbf{RQ3:} Which user personality/archetype modulate the effectiveness of different ad formats?
    \item \textbf{RQ4:} Are there any neural correlates between ad format and cognitive load, attention, and affective arousal (e.g., specific EEG frequency bands and eye-tracking metrics)? Do these signals provide independent, behavior-agnostic validation of user trust and overload beyond conventional self-report methods?
    \item \textbf{RQ5:} How is ad timing (``ad opportunity score'' model) affecting ad conversion?
    \item \textbf{RQ6:} How is ad timing affecting the subjective user experience, including perceived trust in the AI agent?
    \item \textbf{RQ7:} Which user personality/archetype modulate the effectiveness of different ad timing?
    \item \textbf{RQ8:} Are there any neural correlates between ad timing and cognitive load, attention, and affective arousal (e.g., specific EEG frequency bands and eye-tracking metrics)? Do these signals provide independent, behavior-agnostic validation of user trust and overload beyond conventional self-report methods?
    \item \textbf{RQ9:} How is advertising greediness affecting ad conversion?
    \item \textbf{RQ10:} How is advertising greediness affecting the subjective user experience, including perceived trust in the AI agent?
    \item \textbf{RQ:} Which user personality/archetype modulate the effectiveness of different ad greediness?
    \item \textbf{RQ11:} Are there any neural correlates between advertising greediness and cognitive load, attention, and affective arousal (e.g., specific EEG frequency bands and eye-tracking metrics)? Do these signals provide independent, behavior-agnostic validation of user trust and overload beyond conventional self-report methods?
    \item \textbf{RQ12:} In advertisement opportunity modelling are neurophysiological signals and others (opinion of user) or personality a significant feature? (ablation study) 
\end{itemize}

\section{Research Variables}

\begin{itemize}
    \item \textbf{Advertising type/format (categorical; five levels)}:
    \begin{itemize}
        \item \textit{Classical UI advertising}: traditional ad units surrounding the LLM, such as banners, sidebars, sponsored cards, or search-style placements.
        \item \textit{Contextual in-chat advertising}: ads inserted into the conversational flow based on the user’s query or dialogue context, as in \textbf{Bing}-style sponsored placements, and are disclosed.
        \item \textit{Sponsored conversational suggestions}: sponsored follow-up questions or recommendations that appear as part of the interaction, but remain clearly labeled and distinct from the model’s core answer, as in \textbf{Perplexity’s} sponsored search-style recommendations.
        \item \textit{Separated answer with adjacent ads}: the model produces the main response, while ads or recommendations are displayed in a separate area beside or below the answer, as in \textbf{OpenAI}-style side/recommendation modules.
        \item \textit{Implicit or sneaky persuasion}: advertising is blended into the generated response itself with weak separation and no disclosing of it, making the promotional intent less explicit.
    \end{itemize}
        
    \item \textbf{Advertising greediness (categorical; 4 levels)} 
    \begin{itemize}
        \item \textbf{AI‑1 Factual}: Only product‑factual information is provided; no explicit recommendations or CTAs.
        \item \textbf{AI‑2 Low‑intensity}: At most one hedged recommendation per session, triggered only on explicit user preference signals (e.g., preference expressions, comparisons).
        \item \textbf{AI‑3 Medium‑intensity}: Up to 3 recommendations per session, inserted when product mentions are semantically adjacent to the user’s query or intent.
        \item \textbf{AI‑4 High‑intensity}: Recommendations can appear in many turns, may include stronger CTAs, and are less constrained by explicit user intent signals.
    \end{itemize}

    \item \textbf{User personality based on the Big Five (OCEAN) model (five ordinal? dimensions)}: 
    \begin{itemize}
        \item Openness
        \item Conscientiousness
        \item Extraversion
        \item Agreeableness
        \item Neuroticism
    \end{itemize}
    
    \item \textbf{Ad timing $t^*$ (ratio; continuous)}: The turn at which the first recommendation is inserted. This variable is highly related with the modelling of opportunities (which isnt a variable)
\end{itemize}

\section{Research Questions $\times$ Research Variables}

\begin{center}
\begin{tabular}{
    >{\centering\arraybackslash}m{1.2cm}
    |>{\centering\arraybackslash}m{1.2cm}
    |>{\centering\arraybackslash}m{1.2cm}
    |>{\centering\arraybackslash}m{1.2cm}
    |>{\centering\arraybackslash}m{1.2cm}|
}
\multicolumn{1}{c}{} 
& \textbf{V1} & \textbf{V2} & \textbf{V3} & \textbf{V4} \\ \cline{2-5}

\textbf{RQ1}  & $\bullet$ & & & \\ \hline
\textbf{RQ2}  & $\bullet$ & & & \\ \hline
\textbf{RQ3}  & $\bullet$ & & $\bullet$ & \\ \hline
\textbf{RQ4}  & $\bullet$ & & & \\ \hline
\textbf{RQ5}  & & & & $\bullet$ \\ \hline
\textbf{RQ6}  & & & & $\bullet$ \\ \hline
\textbf{RQ7}  & & & $\bullet$ & $\bullet$ \\ \hline
\textbf{RQ8}  & & & & $\bullet$ \\ \hline
\textbf{RQ9}  & & $\bullet$ & & \\ \hline
\textbf{RQ10} & & $\bullet$ & & \\ \hline
\textbf{RQ11} & & $\bullet$ & & \\ \hline
\textbf{RQ12} & & & & \\ \hline
\end{tabular}
\end{center}

\section{Experimentation}

\subsection{Crowd-sourcing Study}

    \begin{itemize}
        \item RQ1
        \item RQ2
        \item RQ3
        \item RQ5
        \item RQ6
        \item RQ7
        \item RQ9
        \item RQ10
    \end{itemize}
    
\subsection{User Study}

    \begin{itemize}
        \item RQ4
        \item RQ8
        \item RQ11
        \item RQ12
    \end{itemize}
    
\section{Contributions}

\begin{enumerate}
    \item Define a theoretical foundation for advertisement in LLM's so big companies or others with data and resources can fucking do something about it 
    \item How people react to advertisement types?
    \item How can we model ad probability (and thus see the importance of each variable (turn to recommendation, eeg signals, eye data,
\end{enumerate}

% --------------------------------------------------
\section{Problem Formalization (Draft)}

\subsection{New Definitions}

\begin{itemize}
    \item  Attention Shift is the causal change in the distribution of focus over latent intents induced by an intervention.
    \\ 
            
        \textbf{Latent Space:}
        \[
        \mathcal{Z} = \text{latent intent / concept space}
        \]
        
        \textbf{Conversation Sequence:}
        \[
        C_{\leq t} = (m_1, m_2, \dots, m_t)
        \]
        where each \(m_i\) is an message (user or assistant message).
        
        \textbf{Attention State:}
        \[
        A_t = P(Z \mid C_{\leq t})
        \]
        
        \textbf{Pre/Post Intervention Attention:}
        \[
        A_{\text{pre}} = P(Z \mid C_{\text{pre}})
        \]
        \[
        A_{\text{post}} = P(Z \mid C_{\text{post}})
        \]
        
        \textbf{Attention Shift (scalar via divergence):}
        \[
        \Delta_{\text{attn}} = D\big(A_{\text{post}} \,\|\, A_{\text{pre}}\big)
        \]
    \item Idk

\end{itemize}

 
\subsection{Entities}
 
\begin{itemize}
    \item User: $u \in \mathcal{U}$, characterized by archetype $\phi(u) \in \{\text{Maximizer, Satisficer, Explorer}\}$
    \item Items: $i \in \mathcal{I}$ (Amazon Electronics catalog)
    \item Retrieval corpus: $\mathcal{D}$ — product titles, descriptions, and aggregated reviews
    \item Conversation history at turn $t$: $c_t = (q_1, a_1, \ldots, q_t)$, bounded at $T = 10$
    \item LLM: $f_\theta$ — prompt-engineered; not fine-tuned
    \item Advertising intensity: $\alpha \in \{\text{AI-1, AI-2, AI-3, AI-4}\}$
    \item Advertising prompt template: $\pi_\alpha$ — system-level instruction encoding intensity level $\alpha$
\end{itemize}
 
\subsection{Interaction Process}
 
At each turn $t$ the system performs two steps:
 
\paragraph{Step 1 — Retrieval.} Given the conversation history $c_t$ and a retrieval strategy RS, the system fetches the top-$k$ candidate items from the corpus:
 
\begin{equation}
    \hat{\mathcal{I}}_t = R(c_t, \mathcal{D}, \text{RS})
\end{equation}
 
The scoring function $s(i, c_t)$ inside $R$ depends on the chosen Recommendation Strategy (Impression-Aware vs.\ Catalog-Aware) — to be formalized once RS is finalized.
 
\paragraph{Step 2 — Generation.} The LLM generates a response conditioned on the conversation, the retrieved items, and the advertising intensity prompt:
 
\begin{equation}
    a_t = f_\theta\!\left(c_t,\; \hat{\mathcal{I}}_t,\; \pi_\alpha\right)
\end{equation}
 
\subsection{Objective (to refine)}
 
The core tension of the system is a \textbf{constrained optimisation}: maximise short-term business reward (CTR) without sacrificing long-term user relationship (trust and retention):
 
\begin{equation}
    \max_{\alpha,\, \text{RS},\, t^*}\; \mathbb{E}_{u}\!\left[\text{CTR}(u, a_t)\right]
    \quad \text{s.t.} \quad
    \mathbb{E}_{u}\!\left[\text{Trust}(u, c_T)\right] \geq \tau
\end{equation}
 
where $\tau$ is an acceptable trust floor (operationalization TBD) and $t^*$ is the turn at which the first recommendation is injected. The decision variables are therefore the triple $(\alpha, \text{RS}, t^*)$.
 
\medskip\noindent
Candidate reward signals (to be narrowed down):
\begin{itemize}
    \item \textbf{CTR}: click or explicit purchase intent signal collected post-session
    \item \textbf{Trust}: self-reported trust scale (TBD) and/or implicit behavioural proxies (e.g.\ conversation abandonment rate)
    \item \textbf{Retention}: return rate across sessions — harder to measure in a lab study; may be simulated
    \item \textbf{Engagement}: conversation length, number of follow-up queries
\end{itemize}
 

\newpage 

% --------------------------------------------------

\begin{thebibliography}{9}

\bibitem{xu2026ad}
Shengwei Xu, Zhaohua Chen, Xiaotie Deng, Zhiyi Huang, and Grant Schoenebeck.
\textit{Ad Insertion in LLM-Generated Responses}.
arXiv preprint arXiv:2601.19435, 2026.
doi: 10.48550/arXiv.2601.19435.

\bibitem{heineking2026detecting}
Sebastian Heineking, Wilhelm Pertsch, Ines Zelch, Janek Bevendorff, Benno Stein, Matthias Hagen, and Martin Potthast.
\textit{Detecting RAG Advertisements Across Advertising Styles}.
arXiv preprint arXiv:2603.04925, 2026.
doi: 10.48550/arXiv.2603.04925.

\bibitem{zhao2025exploring}
Xiaoyan Zhao, Yang Deng, Wenjie Wang, Hongzhan Lin, Hong Cheng, Rui Zhang, See-Kiong Ng, and Tat-Seng Chua.
\textit{Exploring the Impact of Personality Traits on Conversational Recommender Systems: A Simulation with Large Language Models}.
arXiv preprint arXiv:2504.12313, 2025.
doi: 10.48550/arXiv.2504.12313.

\bibitem{feizi2023online}
Soheil Feizi, MohammadTaghi Hajiaghayi, Keivan Rezaei, and Suho Shin.
\textit{Online Advertisements with LLMs: Opportunities and Challenges}.
arXiv preprint arXiv:2311.07601, 2023.
doi: 10.48550/arXiv.2311.07601.

\bibitem{martina2025distillrecdial}
Alessandro Francesco Maria Martina, Alessandro Petruzzelli, Cataldo Musto, Marco de Gemmis, Pasquale Lops, and Giovanni Semeraro.
\textit{DistillRecDial: A Knowledge-Distilled Dataset Capturing User Diversity in Conversational Recommendation}.
In \textit{Proceedings of the 19th ACM Conference on Recommender Systems}, 2025.
ACM.
doi: 10.1145/3705328.3748161.

\end{thebibliography}

\end{thebibliography}

\end{document}